\FloatBarrier
\section{定量申明溯源表}\label{app:provenance}

本表为论文每一处定量申明建立一行溯源：申明内容与所在位置（申明下小字）、生成它的数据文件或脚本、列过滤与聚合口径、样本量与核对状态。全部 115 条均已对照仓库源数据复核。路径相对于仓库根目录：聚合数据位于 \texttt{data/aggregated/}，原始记录位于 \texttt{data/raw/}，脚本位于 \texttt{code/}；派生行（无独立文件）注明其来源行号。论文实验编号与仓库代号、脚本、数据目录的对照见仓库根目录 \texttt{EXPERIMENT\_REGISTRY.md}。

{\footnotesize
\setlength{\LTleft}{0pt}\setlength{\LTright}{0pt}\setlength{\tabcolsep}{3pt}
\begin{longtable}{@{}p{0.9cm}p{6.4cm}p{4.2cm}p{3.0cm}p{1.2cm}@{}}
\caption{定量申明溯源表（共 115 条，全部已核）。}\label{tab:provenance}\\
\toprule
编号 & 申明（下起小字为正文位置） & 数据来源 & 口径（过滤；聚合；$n$） & 状态\\
\midrule
\endfirsthead
\multicolumn{5}{@{}l}{\footnotesize 续表~\ref{tab:provenance}}\\
\toprule
编号 & 申明（下起小字为正文位置） & 数据来源 & 口径（过滤；聚合；$n$） & 状态\\
\midrule
\endhead
\midrule\multicolumn{5}{r@{}}{\footnotesize 接下页}\\
\endfoot
\bottomrule
\endlastfoot
C01 & 实验 4.2 多起点找到的改善 N$\le$512 \newline {\scriptsize §1 摘要/贡献/§4.4} & {\scriptsize\texttt{aggregated/\hspace{0pt}v7\_\hspace{0pt}v2\_\hspace{0pt}exp\_\hspace{0pt}V2\_\hspace{0pt}F1d\_\hspace{0pt}real\_\hspace{0pt}jacobian/\hspace{0pt}f1d\_\hspace{0pt}real\_\hspace{0pt}jacobian.csv}} & nu=\allowbreak .005 且 N$\in${64;\allowbreak 128;\allowbreak 256;\allowbreak 512};\allowbreak 1-ratio\_\allowbreak r;\allowbreak  范围;\allowbreak  n=\allowbreak 4 格 & 已核\\
C02 & 实验 4.2 改善比 r=$\kappa$*\_diag/$\kappa$(R) \newline {\scriptsize §1 摘要/贡献/§4.3/§4.4} & {\scriptsize\texttt{aggregated/\hspace{0pt}v7\_\hspace{0pt}v2\_\hspace{0pt}exp\_\hspace{0pt}V2\_\hspace{0pt}F1d\_\hspace{0pt}real\_\hspace{0pt}jacobian/\hspace{0pt}f1d\_\hspace{0pt}real\_\hspace{0pt}jacobian.csv}} & nu=\allowbreak .005 且 N$\in${64;\allowbreak 128;\allowbreak 256;\allowbreak 512};\allowbreak ratio\_\allowbreak r;\allowbreak  范围;\allowbreak  n=\allowbreak 4 格 & 已核\\
C03 & 实验 4.2 N=1024 时 $\kappa$(K\_r)$\gtrsim$1e13 \newline {\scriptsize §1 摘要/§4.4} & {\scriptsize\texttt{aggregated/\hspace{0pt}v7\_\hspace{0pt}v2\_\hspace{0pt}exp\_\hspace{0pt}V2\_\hspace{0pt}F1d\_\hspace{0pt}real\_\hspace{0pt}jacobian/\hspace{0pt}f1d\_\hspace{0pt}real\_\hspace{0pt}jacobian.csv}} & N=\allowbreak 1024 行 kappa\_\allowbreak R;\allowbreak  原值;\allowbreak  n=\allowbreak 1 & 已核\\
C04 & uniform 可精修盆地命中率 \newline {\scriptsize §1 摘要/§5} & {\scriptsize\texttt{aggregated/\hspace{0pt}v7\_\hspace{0pt}v1\_\hspace{0pt}exp\_\hspace{0pt}G1b\_\hspace{0pt}pentry/\hspace{0pt}g1b\_\hspace{0pt}pentry\_\hspace{0pt}agg.csv}} & scope=\allowbreak pooled;\allowbreak family=\allowbreak uniform;\allowbreak p\_\allowbreak entry(\allowbreak =\allowbreak good\_\allowbreak k/\allowbreak n)\allowbreak ;\allowbreak  比例+Wilson95;\allowbreak  n=\allowbreak 80 & 已核\\
C05 & 加权族命中率 rational/down\_inv/down\_lin \newline {\scriptsize §5} & {\scriptsize\texttt{aggregated/\hspace{0pt}v7\_\hspace{0pt}v1\_\hspace{0pt}exp\_\hspace{0pt}G1b\_\hspace{0pt}pentry/\hspace{0pt}g1b\_\hspace{0pt}pentry\_\hspace{0pt}agg.csv}} & scope=\allowbreak pooled;\allowbreak family 各行 p\_\allowbreak entry;\allowbreak  比例;\allowbreak  n=\allowbreak 80/\allowbreak 族 & 已核\\
C06 & McNemar 净救回 down\_inv/down\_lin/rational \newline {\scriptsize §5} & {\scriptsize\texttt{aggregated/\hspace{0pt}v7\_\hspace{0pt}v1\_\hspace{0pt}exp\_\hspace{0pt}G1b\_\hspace{0pt}pentry/\hspace{0pt}g1b\_\hspace{0pt}flip.csv}} & net 列;\allowbreak  配对计数;\allowbreak  n=\allowbreak 80 & 已核\\
C07 & McNemar 精确双侧 p down\_lin/down\_inv/rational \newline {\scriptsize §5} & {\scriptsize\texttt{aggregated/\hspace{0pt}v7\_\hspace{0pt}v1\_\hspace{0pt}exp\_\hspace{0pt}G1b\_\hspace{0pt}pentry/\hspace{0pt}g1b\_\hspace{0pt}flip.csv}} & 由 (\allowbreak rescued\_\allowbreak BtoG;\allowbreak lost\_\allowbreak GtoB)\allowbreak  精确二项重算;\allowbreak  双侧精确;\allowbreak  n=\allowbreak 80 & 已核\\
C08 & Holm 校正后仅 down\_lin 显著 p \newline {\scriptsize §5} & {\scriptsize 由 C07 三族 p 值 Holm 重算} & 最小 p$\times$3;\allowbreak  Holm;\allowbreak  n=\allowbreak 3 比较 & 已核\\
C09 & 小粘性 $\nu$=.001 配对差 \newline {\scriptsize §5 表注} & {\scriptsize\texttt{aggregated/\hspace{0pt}v7\_\hspace{0pt}v1\_\hspace{0pt}exp\_\hspace{0pt}G1b\_\hspace{0pt}pentry/\hspace{0pt}g1b\_\hspace{0pt}pentry\_\hspace{0pt}agg.csv}} & scope=\allowbreak by\_\allowbreak nu;\allowbreak nu=\allowbreak .001;\allowbreak uniform-各族 p\_\allowbreak entry;\allowbreak  配对差;\allowbreak  n=\allowbreak 20/\allowbreak 族 & 已核\\
C10 & 过渡粘性 $\nu$=.01 方向不一致 \newline {\scriptsize §5 表注} & {\scriptsize\texttt{aggregated/\hspace{0pt}v7\_\hspace{0pt}v1\_\hspace{0pt}exp\_\hspace{0pt}G1b\_\hspace{0pt}pentry/\hspace{0pt}g1b\_\hspace{0pt}pentry\_\hspace{0pt}agg.csv}} & scope=\allowbreak by\_\allowbreak nu;\allowbreak nu=\allowbreak .01;\allowbreak uniform-各族 p\_\allowbreak entry;\allowbreak  配对差;\allowbreak  n=\allowbreak 20/\allowbreak 族 & 已核\\
C11 & 实验 6.1 proxy(g\_et) AUC \newline {\scriptsize §6} & {\scriptsize\texttt{aggregated/\hspace{0pt}v7\_\hspace{0pt}v1\_\hspace{0pt}exp\_\hspace{0pt}G2\_\hspace{0pt}diagnose/\hspace{0pt}g2\_\hspace{0pt}feature\_\hspace{0pt}auc.csv}} & feature=\allowbreak ge\_\allowbreak t;\allowbreak auc;\allowbreak  原值 & 已核\\
C12 & 实验 6.1 oracle AUC \newline {\scriptsize §6} & {\scriptsize\texttt{aggregated/\hspace{0pt}v7\_\hspace{0pt}v1\_\hspace{0pt}exp\_\hspace{0pt}G2\_\hspace{0pt}diagnose/\hspace{0pt}g2\_\hspace{0pt}samples.csv}} & oracle 特征重算;\allowbreak  重算 & 已核\\
C13 & 实验 6.1 灵敏度/特异度 \newline {\scriptsize §6} & {\scriptsize\texttt{aggregated/\hspace{0pt}v7\_\hspace{0pt}v1\_\hspace{0pt}exp\_\hspace{0pt}G2\_\hspace{0pt}diagnose/\hspace{0pt}g2\_\hspace{0pt}samples.csv}} & 阈值口径重算;\allowbreak  重算 & 已核\\
C14 & 实验 5.6 全量 runs 数 \newline {\scriptsize §5} & {\scriptsize\texttt{aggregated/\hspace{0pt}v7\_\hspace{0pt}v1\_\hspace{0pt}exp\_\hspace{0pt}E1b\_\hspace{0pt}smooth/\hspace{0pt}e1b\_\hspace{0pt}runs.csv}} & 数据行数;\allowbreak  计数;\allowbreak  n=\allowbreak 1584 & 已核\\
C15 & 宽度64基准组中位 $\beta$=1$\to$8 \newline {\scriptsize §1 贡献} & {\scriptsize\texttt{aggregated/\hspace{0pt}v7\_\hspace{0pt}v1\_\hspace{0pt}exp\_\hspace{0pt}R1\_\hspace{0pt}arch/\hspace{0pt}r1\_\hspace{0pt}runs.csv}} & axis=\allowbreak width;\allowbreak level=\allowbreak 64;\allowbreak nu=\allowbreak .005;\allowbreak family=\allowbreak rational;\allowbreak $\beta$=\allowbreak 1.0 vs 8.0;\allowbreak L2 列;\allowbreak  组中位;\allowbreak  n=\allowbreak 5/\allowbreak 组 & 已核\\
C16 & 跨粘性表 9 行数字 \newline {\scriptsize §4.4} & {\scriptsize\texttt{aggregated/\hspace{0pt}exp\_\hspace{0pt}V2\_\hspace{0pt}P1\_\hspace{0pt}constrained/\hspace{0pt}p1\_\hspace{0pt}constrained.csv}} & $\nu$$\in${.005;\allowbreak .01;\allowbreak .05}$\times$N$\in${128;\allowbreak 256;\allowbreak 512};\allowbreak  原值;\allowbreak  n=\allowbreak 9 & 已核\\
C17 & Adam 瞬态每步相对漂移/n=500 累计 $\delta$\_P \newline {\scriptsize §4.2 注记} & {\scriptsize\texttt{aggregated/\hspace{0pt}v7\_\hspace{0pt}v2\_\hspace{0pt}exp\_\hspace{0pt}V2\_\hspace{0pt}lem41\_\hspace{0pt}drift\_\hspace{0pt}verify/\hspace{0pt}lem41\_\hspace{0pt}drift\_\hspace{0pt}nu0.005.csv}} & deltaP\_\allowbreak start\_\allowbreak n500;\allowbreak 均值与 均值/\allowbreak 500;\allowbreak  均值;\allowbreak  n=\allowbreak 3 & 已核\\
C18 & Adam 稳态(step$\ge$1000)累计漂移中位 n=50/100/500 \newline {\scriptsize §4.2 注记} & {\scriptsize\texttt{aggregated/\hspace{0pt}v7\_\hspace{0pt}v2\_\hspace{0pt}exp\_\hspace{0pt}V2\_\hspace{0pt}lem41\_\hspace{0pt}drift\_\hspace{0pt}verify/\hspace{0pt}lem41\_\hspace{0pt}drift\_\hspace{0pt}nu0.005.csv}} & deltaP\_\allowbreak mid\_\allowbreak n50/\allowbreak 100/\allowbreak 500;\allowbreak  中位;\allowbreak  n=\allowbreak 3 & 已核\\
C19 & 稳态每步平均漂移 \newline {\scriptsize §4.2 注记} & {\scriptsize （派生自 C18）} & 中位/\allowbreak n 导出;\allowbreak  导出;\allowbreak  n=\allowbreak 3 & 已核\\
C20 & 平台项/白化梯度比值 中位与 IQR \newline {\scriptsize §4.2 注记} & {\scriptsize\texttt{aggregated/\hspace{0pt}v7\_\hspace{0pt}v2\_\hspace{0pt}exp\_\hspace{0pt}V2\_\hspace{0pt}s3\_\hspace{0pt}platform\_\hspace{0pt}ratio/\hspace{0pt}s3\_\hspace{0pt}ratios.csv}} & ratio 列(\allowbreak 汇总见 raw 下 s3\_\allowbreak summary.json)\allowbreak ;\allowbreak  中位+IQR;\allowbreak  n=\allowbreak 150=\allowbreak 3$\times$50 & 已核\\
C21 & N=2 闭式例 $\kappa$(R) 与 t*(a;b) \newline {\scriptsize §4.3} & {\scriptsize 解析重算(无 CSV)} & U=\allowbreak 旋转30$^\circ$;\allowbreak K=\allowbreak U diag(\allowbreak 1;\allowbreak 100)\allowbreak  U\textasciicircum{}T;\allowbreak  重算 & 已核\\
C22 & t=1/a 对照 $\kappa$ 与 (tr)$^2$/det \newline {\scriptsize §4.3 附录} & {\scriptsize 解析重算(无 CSV)} & 2$\times$2 恒等式 $\kappa$+$\kappa$$^{-1}$+2=\allowbreak (\allowbreak tr)\allowbreak $^2$/\allowbreak det;\allowbreak  重算 & 已核\\
C23 & 实验 4.1 等相关阵 42 组 r=1 最大偏离 \newline {\scriptsize §4.3} & {\scriptsize\texttt{aggregated/\hspace{0pt}exp\_\hspace{0pt}V2\_\hspace{0pt}P8\_\hspace{0pt}synthetic\_\hspace{0pt}equicorr/\hspace{0pt}p8\_\hspace{0pt}synthetic\_\hspace{0pt}equicorr.csv}} & r\_\allowbreak ratio 列;\allowbreak max|r-1|;\allowbreak  极值;\allowbreak  n=\allowbreak 42=\allowbreak 6N$\times$7$\rho$ & 已核\\
C24 & 实验 4.1 起点数配置 N$\le$32/64/128 \newline {\scriptsize §4.3} & {\scriptsize\texttt{aggregated/\hspace{0pt}exp\_\hspace{0pt}V2\_\hspace{0pt}P8\_\hspace{0pt}synthetic\_\hspace{0pt}equicorr/\hspace{0pt}p8\_\hspace{0pt}synthetic\_\hspace{0pt}equicorr.csv}} & nstart 列按 N 分档;\allowbreak  逐档;\allowbreak  n=\allowbreak 42 & 已核\\
C25 & 实验 4.1 带状阵 r(30 起点;PSD) \newline {\scriptsize §4.3} & {\scriptsize\texttt{aggregated/\hspace{0pt}exp\_\hspace{0pt}V2\_\hspace{0pt}P8\_\hspace{0pt}synthetic\_\hspace{0pt}equicorr/\hspace{0pt}p8\_\hspace{0pt}banded\_\hspace{0pt}corr.csv}} & r\_\allowbreak ratio 列两行;\allowbreak  原值;\allowbreak  n=\allowbreak 2 & 已核\\
C26 & 实验 4.2 三元组表 $\nu$=.005 全部 24 数 \newline {\scriptsize §4.4 表4.1} & {\scriptsize\texttt{aggregated/\hspace{0pt}exp\_\hspace{0pt}V2\_\hspace{0pt}P1\_\hspace{0pt}constrained/\hspace{0pt}p1\_\hspace{0pt}constrained.csv}} & nu=\allowbreak .005 四行(\allowbreak N=\allowbreak 64--512)\allowbreak ;\allowbreak kappa\_\allowbreak uniform/\allowbreak kappa\_\allowbreak jacobi/\allowbreak kappa\_\allowbreak unconstr/\allowbreak r\_\allowbreak unconstr/\allowbreak kappa\_\allowbreak rational\_\allowbreak best/\allowbreak rational\_\allowbreak best\_\allowbreak beta;\allowbreak  原值;\allowbreak  n=\allowbreak 4 & 已核\\
C27 & uniform/Jacobi 倍数 N=64--128 与 N=256--512 \newline {\scriptsize §4.4} & {\scriptsize\texttt{aggregated/\hspace{0pt}exp\_\hspace{0pt}V2\_\hspace{0pt}P1\_\hspace{0pt}constrained/\hspace{0pt}p1\_\hspace{0pt}constrained.csv}} & nu=\allowbreak .005;\allowbreak uniform\_\allowbreak vs\_\allowbreak jacobi 列;\allowbreak  原值;\allowbreak  n=\allowbreak 4 & 已核\\
C28 & 线性 gap 口径下 Jacobi 填补比例 \newline {\scriptsize §4.4} & {\scriptsize （派生自 C26）} & 1-(\allowbreak $\kappa$*diag-$\kappa$(\allowbreak R)\allowbreak )\allowbreak /\allowbreak (\allowbreak $\kappa$(\allowbreak K\_\allowbreak r)\allowbreak -$\kappa$(\allowbreak R)\allowbreak )\allowbreak ;\allowbreak  导出;\allowbreak  n=\allowbreak 4 & 已核\\
C29 & rational 族内最优：N=64 $\beta$*=5 为 Jacobi 32 倍；N$\ge$128 $\beta$*=1 且单调恶化 \newline {\scriptsize §4.4} & {\scriptsize\texttt{aggregated/\hspace{0pt}exp\_\hspace{0pt}V2\_\hspace{0pt}P1\_\hspace{0pt}constrained/\hspace{0pt}p1\_\hspace{0pt}constrained.csv}} & rational\_\allowbreak best\_\allowbreak beta/\allowbreak r\_\allowbreak rational/\allowbreak kappa\_\allowbreak beta\_\allowbreak 1..100 列;\allowbreak  原值+逐列;\allowbreak  n=\allowbreak 4 & 已核\\
C30 & 有效秩 PR 饱和（N$\ge$256） \newline {\scriptsize §4.4} & {\scriptsize\texttt{aggregated/\hspace{0pt}v7\_\hspace{0pt}v2\_\hspace{0pt}exp\_\hspace{0pt}V2\_\hspace{0pt}F1d\_\hspace{0pt}real\_\hspace{0pt}jacobian/\hspace{0pt}f1d\_\hspace{0pt}real\_\hspace{0pt}jacobian.csv}} & nu=\allowbreak .005 且 N$\in${256;\allowbreak 512;\allowbreak 1024;\allowbreak 2048};\allowbreak eff\_\allowbreak rank\_\allowbreak PR;\allowbreak  原值;\allowbreak  n=\allowbreak 4 & 已核\\
C31 & mpmath 高精度验证 $\lambda$min 相对差异 \newline {\scriptsize §4.4} & {\scriptsize aggregated/\hspace{0pt}v7\_\hspace{0pt}v2\_\hspace{0pt}exp\_\hspace{0pt}V2\_\hspace{0pt}F1d\_\hspace{0pt}N1024\_\hspace{0pt}verify/\hspace{0pt}f1d\_\hspace{0pt}N1024\_\hspace{0pt}verify.csv 与 f1d\_\hspace{0pt}N1024\_\hspace{0pt}verify\_\hspace{0pt}lmin\_\hspace{0pt}allseeds.csv} & lmin\_\allowbreak rel\_\allowbreak diff 列;\allowbreak  原值+补跑;\allowbreak  n=\allowbreak 8 种子 & 已核\\
C32 & 实验 4.3 减权族恶化比值：默认裁剪/小粘性/裁剪关/$\nu$=.01 down\_inv \newline {\scriptsize §4.4/§1} & {\scriptsize\texttt{aggregated/\hspace{0pt}v7\_\hspace{0pt}v1\_\hspace{0pt}exp\_\hspace{0pt}Gstab\_\hspace{0pt}grad/\hspace{0pt}gstab\_\hspace{0pt}agg.csv}} & stage=\allowbreak E;\allowbreak kappa\_\allowbreak med 族÷uniform；clip=\allowbreak 1 全 $\nu$；clip=\allowbreak 1 $\nu$$\in${.001;\allowbreak .005}；clip=\allowbreak 0 全 $\nu$；clip=\allowbreak 0 $\nu$=\allowbreak .01;\allowbreak  比值范围;\allowbreak  n=\allowbreak 11 种子 & 已核\\
C33 & 实验 4.3 g\_peak 与 $\lambda$max 排序 rational$\gg$uniform>down\_lin>down\_inv \newline {\scriptsize §4.4 图题} & {\scriptsize\texttt{aggregated/\hspace{0pt}v7\_\hspace{0pt}v1\_\hspace{0pt}exp\_\hspace{0pt}Gstab\_\hspace{0pt}grad/\hspace{0pt}gstab\_\hspace{0pt}agg.csv}} & stage=\allowbreak E;\allowbreak g\_\allowbreak peak 与 lmax\_\allowbreak med 逐 (\allowbreak $\nu$;\allowbreak clip)\allowbreak  格排序;\allowbreak  排序核验;\allowbreak  n=\allowbreak 6 格 & 已核\\
C34 & 实验 4.4 训练全程 $\kappa$ 比中位且全程>1 \newline {\scriptsize §4.4} & {\scriptsize\texttt{aggregated/\hspace{0pt}v7\_\hspace{0pt}v1\_\hspace{0pt}exp\_\hspace{0pt}C2\_\hspace{0pt}converged\_\hspace{0pt}sa/\hspace{0pt}c2\_\hspace{0pt}runs.csv}} & beta>1（5/\allowbreak 15/\allowbreak 40 合并）；按 stage 对 improve\_\allowbreak ratio 取中位;\allowbreak  中位;\allowbreak  n=\allowbreak 48/\allowbreak stage & 已核\\
C35 & 实验 4.4 收敛点增权族恶化/down\_lin 改善/Gate 有理 \newline {\scriptsize §4.4} & {\scriptsize\texttt{aggregated/\hspace{0pt}v7\_\hspace{0pt}v1\_\hspace{0pt}exp\_\hspace{0pt}E3\_\hspace{0pt}family\_\hspace{0pt}sa/\hspace{0pt}e3\_\hspace{0pt}spectral.csv}} & point=\allowbreak conv 或 gate；按 family 对 improve\_\allowbreak ratio 取中位（contrast 合并）;\allowbreak  中位;\allowbreak  n=\allowbreak 18/\allowbreak 族·点 & 已核\\
C36 & 实验 4.2 多起点总数随 N 递减 \newline {\scriptsize §4.4} & {\scriptsize code exp\_\hspace{0pt}V2\_\hspace{0pt}F1d\_\hspace{0pt}real\_\hspace{0pt}jacobian.py 行 152} & nstart=\allowbreak 12/\allowbreak 8/\allowbreak 6 加 1 个 x=\allowbreak 0 起点（opt\_\allowbreak diag\_\allowbreak kappa）;\allowbreak  配置核对 & 已核\\
C37 & 谱端三族 $\lambda$max 与 $\kappa$ 与负曲率维数 \newline {\scriptsize §4.4 补充观察} & {\scriptsize\texttt{raw/\hspace{0pt}v7\_\hspace{0pt}v2/\hspace{0pt}exp\_\hspace{0pt}V2\_\hspace{0pt}lanczos\_\hspace{0pt}convergence/\hspace{0pt}lanczos\_\hspace{0pt}gate\_\hspace{0pt}probe.csv}} & 定点重跑（diag\_\allowbreak lanczos\_\allowbreak convergence.py 同口径；ncv=\allowbreak (\allowbreak 61;\allowbreak 24)\allowbreak  与 (\allowbreak 100;\allowbreak 60)\allowbreak ；双起点）；$\lambda$max 比 9.36:1:1/\allowbreak 9.21；|$\lambda$min| 缩放 9.80 倍与 1/\allowbreak 9.80;\allowbreak  重跑;\allowbreak  n=\allowbreak 9 次估计 & 已核\\
C38 & N=1024 时 $\kappa$ 近双精度极限（共 4 处） \newline {\scriptsize §1/§4.4} & {\scriptsize\texttt{aggregated/\hspace{0pt}v7\_\hspace{0pt}v2\_\hspace{0pt}exp\_\hspace{0pt}V2\_\hspace{0pt}F1d\_\hspace{0pt}largeN\_\hspace{0pt}fix/\hspace{0pt}f1d\_\hspace{0pt}largeN\_\hspace{0pt}fix.csv}} & kR\_\allowbreak eigh（N=\allowbreak 1024 行）;\allowbreak  原值;\allowbreak  n=\allowbreak 1 & 已核\\
C39 & 实验 4.2 冻结点好盆地 L2 \newline {\scriptsize §4.4} & {\scriptsize\texttt{aggregated/\hspace{0pt}exp\_\hspace{0pt}V2\_\hspace{0pt}P1\_\hspace{0pt}constrained/\hspace{0pt}p1\_\hspace{0pt}constrained.csv}} & gate\_\allowbreak l2（nu=\allowbreak .005 行）;\allowbreak  原值;\allowbreak  n=\allowbreak 4 行同值 & 已核\\
C40 & 实验 5.2 坏盆地 cos $\alpha$\_S（对 d\_B，$\beta$=1） \newline {\scriptsize §5.3 rem:align-emp / §5.5 实验 5.2 段} & {\scriptsize\texttt{aggregated/\hspace{0pt}v7\_\hspace{0pt}v1\_\hspace{0pt}exp\_\hspace{0pt}G3\_\hspace{0pt}c3prime/\hspace{0pt}g3\_\hspace{0pt}probe.csv}} & beta=\allowbreak =\allowbreak 1 且 truth\_\allowbreak good=\allowbreak =\allowbreak 0;\allowbreak cos\_\allowbreak S;\allowbreak  中位+计数+范围;\allowbreak  n=\allowbreak 18 & 已核\\
C41 & 实验 5.2 好盆地 cos $\alpha$\_S（对 d\_*，$\beta$=1） \newline {\scriptsize §5.3 rem:align-emp / §5.5 实验 5.2 段} & {\scriptsize\texttt{aggregated/\hspace{0pt}v7\_\hspace{0pt}v1\_\hspace{0pt}exp\_\hspace{0pt}G3\_\hspace{0pt}c3prime/\hspace{0pt}g3\_\hspace{0pt}probe.csv}} & beta=\allowbreak =\allowbreak 1 且 truth\_\allowbreak good=\allowbreak =\allowbreak 1;\allowbreak cos\_\allowbreak S;\allowbreak  中位+计数+范围;\allowbreak  n=\allowbreak 26 & 已核\\
C42 & C3 内积 $\langle$$\partial$\_$\beta$ g\_S, g\_S$\rangle$ 逐条符号 \newline {\scriptsize §5.3 rem:align-emp/ fig:r1g3 图题} & {\scriptsize\texttt{aggregated/\hspace{0pt}v7\_\hspace{0pt}v1\_\hspace{0pt}exp\_\hspace{0pt}G3\_\hspace{0pt}c3prime/\hspace{0pt}g3\_\hspace{0pt}probe.csv}} & c3\_\allowbreak inner（全表 176 行）;\allowbreak  计数;\allowbreak  n=\allowbreak 176 & 已核\\
C43 & C3 内积中位随 $\beta$ 变化（好盆地） \newline {\scriptsize fig:r1g3 图题} & {\scriptsize\texttt{aggregated/\hspace{0pt}v7\_\hspace{0pt}v1\_\hspace{0pt}exp\_\hspace{0pt}G3\_\hspace{0pt}c3prime/\hspace{0pt}g3\_\hspace{0pt}agg.csv}} & truth\_\allowbreak good=\allowbreak =\allowbreak 1;\allowbreak c3\_\allowbreak inner\_\allowbreak med 按 beta;\allowbreak  原值;\allowbreak  n=\allowbreak 4 档 & 已核\\
C44 & 定号条件达成率（好盆地 cos>$\varepsilon$\_S 占比） \newline {\scriptsize fig:r1g3 图题} & {\scriptsize\texttt{aggregated/\hspace{0pt}v7\_\hspace{0pt}v1\_\hspace{0pt}exp\_\hspace{0pt}G3\_\hspace{0pt}c3prime/\hspace{0pt}g3\_\hspace{0pt}agg.csv}} & truth\_\allowbreak good=\allowbreak =\allowbreak 1;\allowbreak sign\_\allowbreak hold\_\allowbreak rate 按 beta;\allowbreak  原值;\allowbreak  n=\allowbreak 4 档 & 已核\\
C45 & 实验 5.7 关闭裁剪后零爆炸（含 $\nu$=.001 rational） \newline {\scriptsize §5.4 实验 5.7 段/§5.5} & {\scriptsize\texttt{aggregated/\hspace{0pt}v7\_\hspace{0pt}v2\_\hspace{0pt}exp\_\hspace{0pt}V2\_\hspace{0pt}G1c\_\hspace{0pt}clip\_\hspace{0pt}factor/\hspace{0pt}g1c\_\hspace{0pt}clip\_\hspace{0pt}agg.csv}} & clip=\allowbreak =\allowbreak 0;\allowbreak blown\_\allowbreak k;\allowbreak  原值;\allowbreak  n=\allowbreak 8 格(\allowbreak 2$\nu$$\times$4族)\allowbreak  & 已核\\
C46 & $\nu$=.005 裁剪开关对进入率无影响 \newline {\scriptsize §5.4 实验 5.7 段/§5.5} & {\scriptsize\texttt{g1c\_\hspace{0pt}clip\_\hspace{0pt}runs.csv + v7\_\hspace{0pt}v1\_\hspace{0pt}exp\_\hspace{0pt}G1b\_\hspace{0pt}pentry/\hspace{0pt}g1b\_\hspace{0pt}pentry\_\hspace{0pt}runs.csv}} & $\nu$=\allowbreak .005;\allowbreak 按 seed 配对 clip0 vs clip1(\allowbreak =\allowbreak 实验 5.3)\allowbreak ;\allowbreak good 列;\allowbreak  配对 McNemar 精确;\allowbreak  n=\allowbreak 20 & 已核\\
C47 & $\nu$=.001 clip=1 时 uniform 优于两减权族 \newline {\scriptsize §5.4 实验 5.7 段/§5.5} & {\scriptsize\texttt{v7\_\hspace{0pt}v1\_\hspace{0pt}exp\_\hspace{0pt}G1b\_\hspace{0pt}pentry/\hspace{0pt}g1b\_\hspace{0pt}pentry\_\hspace{0pt}runs.csv + g1c\_\hspace{0pt}clip\_\hspace{0pt}runs.csv}} & $\nu$=\allowbreak .001;\allowbreak 族间按 seed 配对;\allowbreak good 列;\allowbreak  配对 McNemar 精确;\allowbreak  n=\allowbreak 20 & 已核\\
C48 & 实验 5.8/5.9 裁剪触发率 \newline {\scriptsize §5.4 实验 5.7 段} & {\scriptsize\texttt{aggregated/\hspace{0pt}v7\_\hspace{0pt}v2\_\hspace{0pt}exp\_\hspace{0pt}V2\_\hspace{0pt}GstabB\_\hspace{0pt}stabilizer/\hspace{0pt}gstabB\_\hspace{0pt}lr\_\hspace{0pt}runs.csv}} & family=\allowbreak =\allowbreak uniform 且 nu=\allowbreak =\allowbreak .005 且 lr=\allowbreak =\allowbreak 1e-3;\allowbreak frac\_\allowbreak clip;\allowbreak  范围+中位;\allowbreak  n=\allowbreak 8 & 已核\\
C49 & 实验 5.8 爆炸率 \newline {\scriptsize §5.4 实验 5.8 段/§5.5} & {\scriptsize\texttt{aggregated/\hspace{0pt}v7\_\hspace{0pt}v2\_\hspace{0pt}exp\_\hspace{0pt}V2\_\hspace{0pt}GstabB\_\hspace{0pt}stabilizer/\hspace{0pt}gstabB\_\hspace{0pt}lr\_\hspace{0pt}runs.csv}} & blown 按 family$\times$lr 聚合;\allowbreak  比例;\allowbreak  n=\allowbreak 8 种子$\times$4 族 & 已核\\
C50 & 实验 5.8 有效 lr 上限中位（L2<.06） \newline {\scriptsize §5.4 实验 5.8 段/§5.5} & {\scriptsize\texttt{aggregated/\hspace{0pt}v7\_\hspace{0pt}v2\_\hspace{0pt}exp\_\hspace{0pt}V2\_\hspace{0pt}GstabB\_\hspace{0pt}stabilizer/\hspace{0pt}gstabB\_\hspace{0pt}lr\_\hspace{0pt}runs.csv}} & blown=\allowbreak =\allowbreak 0;\allowbreak 每(\allowbreak seed,family)\allowbreak 取 L2<.06 的最大 lr;\allowbreak  中位;\allowbreak  n=\allowbreak 8 & 已核\\
C51 & 实验 5.8 上限配对 Wilcoxon \newline {\scriptsize §5.4 实验 5.8 段/§5.5} & {\scriptsize\texttt{aggregated/\hspace{0pt}v7\_\hspace{0pt}v2\_\hspace{0pt}exp\_\hspace{0pt}V2\_\hspace{0pt}GstabB\_\hspace{0pt}stabilizer/\hspace{0pt}gstabB\_\hspace{0pt}lr\_\hspace{0pt}runs.csv}} & 同 C50 的上限值配对;\allowbreak  精确枚举+脚本渐近（无连续性修正）;\allowbreak  n=\allowbreak 8（非零 5/\allowbreak 2/\allowbreak 3） & 已核\\
C52 & 宽松阈值上限中位（L2<.1/.2） \newline {\scriptsize §5.4 实验 5.8 段} & {\scriptsize\texttt{aggregated/\hspace{0pt}v7\_\hspace{0pt}v2\_\hspace{0pt}exp\_\hspace{0pt}V2\_\hspace{0pt}GstabB\_\hspace{0pt}stabilizer/\hspace{0pt}gstabB\_\hspace{0pt}lr\_\hspace{0pt}runs.csv}} & 同 C50 改阈值;\allowbreak  中位;\allowbreak  n=\allowbreak 8 & 已核\\
C53 & 实验 5.9 $\nu$=.005 四族终态 L2 \newline {\scriptsize §5.4 实验 5.9 段/§5.5} & {\scriptsize\texttt{aggregated/\hspace{0pt}v7\_\hspace{0pt}v2\_\hspace{0pt}exp\_\hspace{0pt}V2\_\hspace{0pt}GstabB\_\hspace{0pt}stabilizer/\hspace{0pt}gstabB\_\hspace{0pt}lbfgs\_\hspace{0pt}runs.csv}} & nu=\allowbreak =\allowbreak .005;\allowbreak L2 按 seed 配对;\allowbreak  中位+脚本渐近 Wilcoxon;\allowbreak  n=\allowbreak 8 & 已核\\
C54 & 实验 5.9 $\nu$=.001 四族终态 L2 \newline {\scriptsize §5.4 实验 5.9 段/§5.5} & {\scriptsize\texttt{aggregated/\hspace{0pt}v7\_\hspace{0pt}v2\_\hspace{0pt}exp\_\hspace{0pt}V2\_\hspace{0pt}GstabB\_\hspace{0pt}stabilizer/\hspace{0pt}gstabB\_\hspace{0pt}lbfgs\_\hspace{0pt}runs.csv}} & nu=\allowbreak =\allowbreak .001;\allowbreak L2 按 seed 配对;\allowbreak  中位+脚本渐近 Wilcoxon;\allowbreak  n=\allowbreak 5 & 已核\\
C55 & GC 500 步梯度范数排序与量级 \newline {\scriptsize §5.4 GC 段} & {\scriptsize\texttt{aggregated/\hspace{0pt}v7\_\hspace{0pt}v2\_\hspace{0pt}exp\_\hspace{0pt}V2\_\hspace{0pt}GstabB\_\hspace{0pt}stabilizer/\hspace{0pt}gstabB\_\hspace{0pt}grad\_\hspace{0pt}curve.csv}} & 按(\allowbreak family,step)\allowbreak 取 8 种子中位;\allowbreak g\_\allowbreak pre;\allowbreak  逐步中位范围;\allowbreak  n=\allowbreak 8$\times$11 步 & 已核\\
C56 & 实验 5.7 rational 裁剪前峰值 7e2（$\nu$=.001） \newline {\scriptsize §5.4 实验 5.7 段} & {\scriptsize\texttt{aggregated/\hspace{0pt}v7\_\hspace{0pt}v1\_\hspace{0pt}exp\_\hspace{0pt}Gstab\_\hspace{0pt}grad/\hspace{0pt}gstab\_\hspace{0pt}agg.csv}} & stage=\allowbreak =\allowbreak E 且 nu=\allowbreak =\allowbreak .001 且 clip=\allowbreak =\allowbreak 0 且 family=\allowbreak =\allowbreak rational;\allowbreak g\_\allowbreak peak;\allowbreak  原值（中位）;\allowbreak  n=\allowbreak 11 & 已核\\
C57 & 冻结点自检脚本 的 $\eta$\_gate 对 $\beta$ Spearman 全为 +1 \newline {\scriptsize §5.5 实验 5.1/5.2 段} & {\scriptsize\texttt{aggregated/\hspace{0pt}v7\_\hspace{0pt}v1\_\hspace{0pt}exp\_\hspace{0pt}Vverify/\hspace{0pt}vv\_\hspace{0pt}claim1\_\hspace{0pt}from\_\hspace{0pt}e1.csv}} & 按 nu$\times$beta 取 spearman 列;\allowbreak  Spearman 秩相关;\allowbreak  n=\allowbreak 48 & 已核\\
C58 & 实验 5.1 跨 7 个配置方向不变 \newline {\scriptsize §5.5 实验 5.1/5.2 段} & {\scriptsize\texttt{aggregated/\hspace{0pt}v7\_\hspace{0pt}v1\_\hspace{0pt}exp\_\hspace{0pt}R1\_\hspace{0pt}arch/\hspace{0pt}r1\_\hspace{0pt}invariance.csv}} & width/\allowbreak depth/\allowbreak npde 共 7 配置;\allowbreak eta\_\allowbreak beta\_\allowbreak spearman;\allowbreak  Spearman 秩相关;\allowbreak  n=\allowbreak 7 & 已核\\
C59 & C-Sep 直接测量：中位比与严格版 \newline {\scriptsize §5.5 实验 5.1/5.2 段} & {\scriptsize\texttt{aggregated/\hspace{0pt}exp\_\hspace{0pt}V2\_\hspace{0pt}P1\_\hspace{0pt}constrained/\hspace{0pt}p1\_\hspace{0pt}constrained.csv}} & nu=\allowbreak .005;\allowbreak N=\allowbreak 128/\allowbreak 256/\allowbreak 512;\allowbreak csep\_\allowbreak median/\allowbreak csep\_\allowbreak strict;\allowbreak  原值;\allowbreak  n=\allowbreak 3 & 已核\\
C60 & down\_lin Holm 显著区间 \newline {\scriptsize §5.5 N6} & {\scriptsize\texttt{aggregated/\hspace{0pt}v7\_\hspace{0pt}v1\_\hspace{0pt}exp\_\hspace{0pt}G1b\_\hspace{0pt}pentry/\hspace{0pt}g1b\_\hspace{0pt}pentry\_\hspace{0pt}runs.csv + code/\hspace{0pt}v7\_\hspace{0pt}1d/\hspace{0pt}experiments/\hspace{0pt}V2/\hspace{0pt}exp\_\hspace{0pt}V2\_\hspace{0pt}N1\_\hspace{0pt}N6\_\hspace{0pt}stats.py}} & tau$\in$[.05,.08];\allowbreak phase0\_\allowbreak L2<tau；配对 McNemar+Holm;\allowbreak  精确 McNemar+标准 Holm（累计 max）;\allowbreak  n=\allowbreak 80 & 已核\\
C61 & 低阈值好盆地样本数 \newline {\scriptsize §5.5 N6} & {\scriptsize\texttt{aggregated/\hspace{0pt}v7\_\hspace{0pt}v1\_\hspace{0pt}exp\_\hspace{0pt}G1b\_\hspace{0pt}pentry/\hspace{0pt}g1b\_\hspace{0pt}pentry\_\hspace{0pt}runs.csv}} & tau=\allowbreak .03/\allowbreak .04；按 family 统计 phase0\_\allowbreak L2<tau;\allowbreak  计数;\allowbreak  n=\allowbreak 80/\allowbreak 族 & 已核\\
C62 & 逐种子最优 $\beta$* 各 $\nu$ 中位 \newline {\scriptsize §5.5 fig:e1heat/附录 实验 5.4} & {\scriptsize\texttt{aggregated/\hspace{0pt}v7\_\hspace{0pt}v1\_\hspace{0pt}exp\_\hspace{0pt}E1\_\hspace{0pt}beta\_\hspace{0pt}nu/\hspace{0pt}e1\_\hspace{0pt}runs.csv}} & nu=\allowbreak .001/\allowbreak .003/\allowbreak .005/\allowbreak .01；每 (\allowbreak nu,seed)\allowbreak  取 L2 最小的 beta;\allowbreak  逐种子 argmin 后中位;\allowbreak  n=\allowbreak 44 & 已核\\
C63 & 跨粘度 $\beta$* 中位 \newline {\scriptsize §5.5 实验 5.4/5.5} & {\scriptsize\texttt{aggregated/\hspace{0pt}v7\_\hspace{0pt}v1\_\hspace{0pt}exp\_\hspace{0pt}E1\_\hspace{0pt}beta\_\hspace{0pt}nu/\hspace{0pt}e1\_\hspace{0pt}runs.csv + aggregated/\hspace{0pt}v7\_\hspace{0pt}v1\_\hspace{0pt}exp\_\hspace{0pt}R2\_\hspace{0pt}budget/\hspace{0pt}r2\_\hspace{0pt}runs.csv}} & 实验 5.4 同 C62；实验 5.5 每 (\allowbreak nu,seed)\allowbreak  取 L2 最小的 beta;\allowbreak  逐种子 argmin 后中位;\allowbreak  n=\allowbreak E1 44；R2 33 & 已核\\
C64 & g\_post 分布与括注 \newline {\scriptsize §5.5 实验 5.4} & {\scriptsize\texttt{aggregated/\hspace{0pt}v7\_\hspace{0pt}v1\_\hspace{0pt}exp\_\hspace{0pt}E1\_\hspace{0pt}beta\_\hspace{0pt}nu/\hspace{0pt}e1\_\hspace{0pt}runs.csv}} & 薄激波 4$\nu$$\times$11；g\_\allowbreak post=\allowbreak (\allowbreak L2(\allowbreak beta=\allowbreak 1)\allowbreak -min\_\allowbreak beta L2)\allowbreak /\allowbreak L2(\allowbreak beta=\allowbreak 1)\allowbreak ;\allowbreak  中位/\allowbreak IQR/\allowbreak 范围;\allowbreak  n=\allowbreak 44 & 已核\\
C65 & 刚性端 4--10 倍增益 \newline {\scriptsize §5.5 实验 5.6} & {\scriptsize\texttt{aggregated/\hspace{0pt}v7\_\hspace{0pt}v1\_\hspace{0pt}exp\_\hspace{0pt}E1b\_\hspace{0pt}smooth/\hspace{0pt}e1b\_\hspace{0pt}runs.csv + e1b\_\hspace{0pt}summary.csv}} & feat=\allowbreak MLP;\allowbreak nu=\allowbreak .01/\allowbreak .03；逐 seed 取 l2\_\allowbreak band 或 L2 最小；gain=\allowbreak x(\allowbreak beta=\allowbreak 1)\allowbreak /\allowbreak best-1;\allowbreak  中位;\allowbreak  n=\allowbreak 22（2$\nu$$\times$11） & 已核\\
C66 & 光滑端退化 \newline {\scriptsize §5.5 实验 5.6} & {\scriptsize\texttt{aggregated/\hspace{0pt}v7\_\hspace{0pt}v1\_\hspace{0pt}exp\_\hspace{0pt}E1b\_\hspace{0pt}smooth/\hspace{0pt}e1b\_\hspace{0pt}csep.csv + e1b\_\hspace{0pt}runs.csv}} & nu=\allowbreak .3/\allowbreak .5/\allowbreak 1；frac\_\allowbreak S；L2;\allowbreak  范围+中位;\allowbreak  n=\allowbreak csep 33/\allowbreak 粘度；runs 66/\allowbreak feat/\allowbreak 粘度 & 已核\\
C67 & 固定高带宽在光滑端失配 \newline {\scriptsize §5.5 实验 5.6} & {\scriptsize\texttt{aggregated/\hspace{0pt}v7\_\hspace{0pt}v1\_\hspace{0pt}exp\_\hspace{0pt}E1b\_\hspace{0pt}smooth/\hspace{0pt}e1b\_\hspace{0pt}runs.csv}} & nu=\allowbreak .3/\allowbreak .5/\allowbreak 1；每 (\allowbreak feat,nu,seed)\allowbreak  取 L2 最小后按 feat 中位再取比;\allowbreak  逐 seed best 中位+比值;\allowbreak  n=\allowbreak 11/\allowbreak 格 & 已核\\
C68 & 中位热力图代表性 \newline {\scriptsize §5.5 fig:e1heat} & {\scriptsize\texttt{aggregated/\hspace{0pt}v7\_\hspace{0pt}v1\_\hspace{0pt}exp\_\hspace{0pt}E1\_\hspace{0pt}beta\_\hspace{0pt}nu/\hspace{0pt}e1\_\hspace{0pt}runs.csv}} & 每 nu：格内中位 argmin vs 逐种子 argmin；逐 seed L2-beta Spearman;\allowbreak  对比;\allowbreak  n=\allowbreak 44 & 已核\\
C69 & 实验 5.4 各格离散度 \newline {\scriptsize §5.5 fig:e1heat} & {\scriptsize\texttt{aggregated/\hspace{0pt}v7\_\hspace{0pt}v1\_\hspace{0pt}exp\_\hspace{0pt}E1\_\hspace{0pt}beta\_\hspace{0pt}nu/\hspace{0pt}e1\_\hspace{0pt}runs.csv + e1\_\hspace{0pt}agg.csv}} & 按 nu$\times$beta 统计 L2 q1/\allowbreak q3/\allowbreak min/\allowbreak max/\allowbreak success;\allowbreak  IQR 与成功率;\allowbreak  n=\allowbreak 11/\allowbreak 格 & 已核\\
C70 & 四点玩具探针 AUC \newline {\scriptsize §6.2 (iii)} & {\scriptsize\texttt{aggregated/\hspace{0pt}v7\_\hspace{0pt}v1\_\hspace{0pt}exp\_\hspace{0pt}G2\_\hspace{0pt}diagnose/\hspace{0pt}g2\_\hspace{0pt}toy\_\hspace{0pt}check.csv}} & toy\_\allowbreak fourpoint(\allowbreak n=\allowbreak 400)\allowbreak  合成探针;\allowbreak  AUC;\allowbreak  n=\allowbreak 400 & 已核\\
C71 & 实验 6.1 标签规模 \newline {\scriptsize §6.3} & {\scriptsize\texttt{raw/\hspace{0pt}v7\_\hspace{0pt}v1/\hspace{0pt}exp\_\hspace{0pt}G2\_\hspace{0pt}diagnose/\hspace{0pt}g2\_\hspace{0pt}samples.csv}} & oracle\_\allowbreak good/\allowbreak truth\_\allowbreak good 列;\allowbreak  计数;\allowbreak  n=\allowbreak 80 & 已核\\
C72 & 49/80 与 实验 5.3 命中率同一统计量 \newline {\scriptsize §6.3} & {\scriptsize\texttt{g2\_\hspace{0pt}samples.csv + aggregated/\hspace{0pt}v7\_\hspace{0pt}v1\_\hspace{0pt}exp\_\hspace{0pt}G1b\_\hspace{0pt}pentry/\hspace{0pt}g1b\_\hspace{0pt}pentry\_\hspace{0pt}runs.csv}} & family=\allowbreak =\allowbreak uniform；按 (\allowbreak nu;\allowbreak seed)\allowbreak  合并 oracle\_\allowbreak L2 vs phase0\_\allowbreak L2;\allowbreak  逐配置对比;\allowbreak  n=\allowbreak 80 & 已核\\
C73 & 7 个训练过程量 AUC \newline {\scriptsize §6.3 (i)} & {\scriptsize\texttt{aggregated/\hspace{0pt}v7\_\hspace{0pt}v1\_\hspace{0pt}exp\_\hspace{0pt}G2\_\hspace{0pt}diagnose/\hspace{0pt}g2\_\hspace{0pt}feature\_\hspace{0pt}auc.csv}} & proxy/\allowbreak loss\_\allowbreak cv/\allowbreak loss\_\allowbreak slope/\allowbreak gradnorm\_\allowbreak t/\allowbreak wmean\_\allowbreak t/\allowbreak focus\_\allowbreak t/\allowbreak focus\_\allowbreak cv 七行 auc;\allowbreak  中位+范围;\allowbreak  n=\allowbreak 80 & 已核\\
C74 & r\_resid 两口径 AUC \newline {\scriptsize §6.3 (ii)} & {\scriptsize\texttt{g2\_\hspace{0pt}feature\_\hspace{0pt}auc.csv + raw g2\_\hspace{0pt}samples.csv}} & r\_\allowbreak resid 对 truth\_\allowbreak good/\allowbreak oracle\_\allowbreak good;\allowbreak  AUC;\allowbreak  n=\allowbreak 80 & 已核\\
C75 & 结构量 proxy AUC \newline {\scriptsize §6.3 (iii)} & {\scriptsize\texttt{g2\_\hspace{0pt}feature\_\hspace{0pt}auc.csv}} & peak\_\allowbreak ratio/\allowbreak peak\_\allowbreak in\_\allowbreak band/\allowbreak abs\_\allowbreak peak/\allowbreak ge\_\allowbreak t 四行 auc;\allowbreak  AUC;\allowbreak  n=\allowbreak 80 & 已核\\
C76 & 结构量 oracle 存在性口径 \newline {\scriptsize §6.3 (iii)} & {\scriptsize\texttt{raw g2\_\hspace{0pt}samples.csv}} & 四特征对 oracle\_\allowbreak good 重算;\allowbreak  AUC;\allowbreak  n=\allowbreak 80 & 已核\\
C77 & g\_et LOO 指标与 17 维持平 \newline {\scriptsize §6.3} & {\scriptsize\texttt{g2\_\hspace{0pt}feature\_\hspace{0pt}auc.csv + g2\_\hspace{0pt}group\_\hspace{0pt}loo.csv}} & ge\_\allowbreak t 行 loo 列；ALL\_\allowbreak 全多变量 行;\allowbreak  LOO Youden；混淆矩阵;\allowbreak  n=\allowbreak 80 & 已核\\
C78 & g\_et proxy AUC bootstrap 区间 \newline {\scriptsize §6.3} & {\scriptsize\texttt{raw g2\_\hspace{0pt}samples.csv}} & ge\_\allowbreak t 对 truth\_\allowbreak good 分层重抽样;\allowbreak  百分位 2.5/\allowbreak 97.5;\allowbreak  n=\allowbreak 80 & 已核\\
C79 & g\_et oracle 口径 \newline {\scriptsize §6.3} & {\scriptsize\texttt{raw g2\_\hspace{0pt}samples.csv}} & ge\_\allowbreak t 对 oracle\_\allowbreak good；LOO Youden;\allowbreak  AUC+LOO 混淆;\allowbreak  n=\allowbreak 80 & 已核\\
C80 & 端到端粗略账本 \newline {\scriptsize §6.3} & {\scriptsize\texttt{g2\_\hspace{0pt}samples.csv}} & oracle\_\allowbreak good 比例$\times$proxy 灵敏度;\allowbreak  算术;\allowbreak  n=\allowbreak 80 & 已核\\
C81 & 好坏盆地 L2 自然间隙与 实验 6.2 \newline {\scriptsize §6.3+附录方向安全声明} & {\scriptsize\texttt{g2\_\hspace{0pt}samples.csv + aggregated/\hspace{0pt}v7\_\hspace{0pt}v1\_\hspace{0pt}exp\_\hspace{0pt}Enew3\_\hspace{0pt}tau\_\hspace{0pt}scan/\hspace{0pt}tau\_\hspace{0pt}scan.csv}} & truth\_\allowbreak L2 按 truth\_\allowbreak good 分组；tau 扫描 agree\_\allowbreak pct;\allowbreak  范围+一致率;\allowbreak  n=\allowbreak 80 & 已核\\
C82 & STRUCT 组内其余量区间 \newline {\scriptsize §6.3 fig:g2auc 图题} & {\scriptsize\texttt{g2\_\hspace{0pt}feature\_\hspace{0pt}auc.csv}} & STRUCT 组除三名外五行 auc;\allowbreak  范围;\allowbreak  n=\allowbreak 80 & 已核\\
C83 & K=4 多起点构造 \newline {\scriptsize §5.5 实验 5.3+§6.3} & {\scriptsize\texttt{code/\hspace{0pt}v7\_\hspace{0pt}1d/\hspace{0pt}experiments/\hspace{0pt}v1/\hspace{0pt}exp\_\hspace{0pt}G1b\_\hspace{0pt}pentry.py + exp\_\hspace{0pt}G2\_\hspace{0pt}basin\_\hspace{0pt}diagnose.py}} & seed=\allowbreak base\_\allowbreak seed+k 循环;\allowbreak  代码审查;\allowbreak  n=\allowbreak 4/\allowbreak 配置 & 已核\\
C84 & 正类规模说明 \newline {\scriptsize §6.3} & {\scriptsize\texttt{raw g2\_\hspace{0pt}samples.csv}} & truth\_\allowbreak good 列;\allowbreak  计数;\allowbreak  n=\allowbreak 80 & 已核\\
C85 & 训练总次数 \newline {\scriptsize §7.1 实验 7.1 设置} & {\scriptsize\texttt{aggregated/\hspace{0pt}v7\_\hspace{0pt}v2\_\hspace{0pt}exp\_\hspace{0pt}V2\_\hspace{0pt}S2a\_\hspace{0pt}poisson\_\hspace{0pt}degenerate/\hspace{0pt}poisson\_\hspace{0pt}runs.csv}} & 全表行数;\allowbreak  计数;\allowbreak  n=\allowbreak 462 & 已核\\
C86 & 增权六组配对五组不显著；显著组（单频 $\sigma$=1）中位 1.41e-5$\to$1.11e-4 p=.001 变差倍数中位约 5 \newline {\scriptsize §7.1 (i)} & {\scriptsize\texttt{poisson\_\hspace{0pt}agg.csv + poisson\_\hspace{0pt}runs.csv}} & 增权 $\beta$=\allowbreak 10 vs 等权 按配置分组;\allowbreak  Wilcoxon 配对+中位;\allowbreak  n=\allowbreak 11/\allowbreak 组 & 已核\\
C87 & 多频+普通MLP 逆减权 7.42e-3$\to$2.02e-4 p=.001 11/11 改善 倍数中位约 24 \newline {\scriptsize §7.1 (ii)} & {\scriptsize 同上} & 多频$\times$MLP 组逆减权 vs 等权;\allowbreak  Wilcoxon 配对+中位;\allowbreak  n=\allowbreak 11 & 已核\\
C88 & 单频 $\sigma$=1 逆减权 1.41e-5$\to$4.05e-6 p=.005 9/11 改善 倍数约 3.3 \newline {\scriptsize §7.1 (ii)} & {\scriptsize 同上} & 单频$\times$$\sigma$=\allowbreak 1 组;\allowbreak  Wilcoxon 配对+中位;\allowbreak  n=\allowbreak 11 & 已核\\
C89 & 失配高带宽（多频 $\sigma$=15）逆减权 .108$\to$.113 p=.001 显著变差 \newline {\scriptsize §7.1 (ii)} & {\scriptsize 同上} & 多频$\times$$\sigma$=\allowbreak 15 组;\allowbreak  Wilcoxon 配对+中位;\allowbreak  n=\allowbreak 11 & 已核\\
C90 & 失配高带宽比匹配带宽差 40--1000 倍 \newline {\scriptsize §7.1 (iii)} & {\scriptsize 同上} & $\sigma$=\allowbreak 15 vs $\sigma$=\allowbreak 1 跨组中位比;\allowbreak  范围;\allowbreak  n=\allowbreak 11/\allowbreak 组 & 已核\\
C91 & 三厚度 K=4 选优 11/11 进可精修盆地 精修后中位 1e-5 量级 \newline {\scriptsize §7.2 实验 7.2} & {\scriptsize\texttt{aggregated/\hspace{0pt}v7\_\hspace{0pt}v2\_\hspace{0pt}exp\_\hspace{0pt}V2\_\hspace{0pt}S2d\_\hspace{0pt}internal\_\hspace{0pt}layer/\hspace{0pt}s2d\_\hspace{0pt}runs.csv}} & 按 $\varepsilon$ 分组终态 L2;\allowbreak  中位+计数;\allowbreak  n=\allowbreak 11/\allowbreak 档 & 已核\\
C92 & 分区梯度能量比 $\eta$ 在 7/11 Gate 点下降 4/11 上升 平均 Spearman 约 -0.32 \newline {\scriptsize §7.2 实验 7.2 机制A条件性} & {\scriptsize\texttt{aggregated/\hspace{0pt}v7\_\hspace{0pt}v2\_\hspace{0pt}exp\_\hspace{0pt}V2\_\hspace{0pt}S2d\_\hspace{0pt}beta\_\hspace{0pt}effect/\hspace{0pt}s2d\_\hspace{0pt}beta\_\hspace{0pt}probe.csv}} & Gate 冻结探针 $\beta$$\in${1;\allowbreak 3;\allowbreak 8};\allowbreak  逐点符号+秩相关;\allowbreak  n=\allowbreak 11 Gate 点 & 已核\\
C93 & $\beta$=3 对 $\beta$=1 为 8/11 更好 中位差 -1.06e-5 p=.41 Holm .56；$\beta$=8 为 7/11 p=.28 Holm .56 \newline {\scriptsize §7.2 实验 7.2 配对精修} & {\scriptsize\texttt{aggregated/\hspace{0pt}v7\_\hspace{0pt}v2\_\hspace{0pt}exp\_\hspace{0pt}V2\_\hspace{0pt}S2d\_\hspace{0pt}beta\_\hspace{0pt}effect/\hspace{0pt}s2d\_\hspace{0pt}beta\_\hspace{0pt}refine\_\hspace{0pt}cmp.csv}} & 层区 L2 配对差;\allowbreak  Wilcoxon+Holm;\allowbreak  n=\allowbreak 11 & 已核\\
C94 & 无种子 12/12 塌零 L2$\approx$.999 残差可降至近零 \newline {\scriptsize §7.2 实验 7.3} & {\scriptsize\texttt{aggregated/\hspace{0pt}v7\_\hspace{0pt}v2\_\hspace{0pt}exp\_\hspace{0pt}V2\_\hspace{0pt}S2e\_\hspace{0pt}spike/\hspace{0pt}s2e\_\hspace{0pt}runs.csv}} & 无种子组终态;\allowbreak  计数+中位;\allowbreak  n=\allowbreak 12 & 已核\\
C95 & 监督对照全域归一化 L2 中位 6.2e-4 较塌零低约三个数量级 \newline {\scriptsize §7.2 实验 7.3} & {\scriptsize\texttt{aggregated/\hspace{0pt}v7\_\hspace{0pt}v2\_\hspace{0pt}exp\_\hspace{0pt}V2\_\hspace{0pt}S2e\_\hspace{0pt}supervised/\hspace{0pt}s2e\_\hspace{0pt}supervised.csv}} & 12 独立种子拟合 u*;\allowbreak  中位;\allowbreak  n=\allowbreak 12 & 已核\\
C96 & 拓扑种子：精确 sech 4/4；粗高斯三幅值均 4/4；错位 x0=.15（约 34$\varepsilon$）4/4 失败 \newline {\scriptsize §7.2 实验 7.3} & {\scriptsize\texttt{s2e\_\hspace{0pt}runs.csv}} & 按种子类型分组;\allowbreak  计数;\allowbreak  n=\allowbreak 4/\allowbreak 组 & 已核\\
C97 & 同厚度 Burgers K=8 的 8 起点 3 命中；贴边层 $\sigma$=15/30 共 16/16 零命中 L2=9.6--10.8 \newline {\scriptsize §7.2 实验 7.4} & {\scriptsize\texttt{aggregated/\hspace{0pt}v7\_\hspace{0pt}v2\_\hspace{0pt}exp\_\hspace{0pt}V2\_\hspace{0pt}S2c\_\hspace{0pt}aligned\_\hspace{0pt}twostage/\hspace{0pt}s2c\_\hspace{0pt}runs.csv}} & 对照组与无干预组终态 L2;\allowbreak  计数+范围;\allowbreak  n=\allowbreak 8/\allowbreak 16 & 已核\\
C98 & 拉伸厚层 4/4 成功 薄层 8/8 失败；精确指数提升 4/4（全域 1.6e-3--3.3e-3）粗 logistic 8/8 失败 \newline {\scriptsize §7.2 实验 7.4} & {\scriptsize s2c\_\hspace{0pt}runs.csv + diag\_\hspace{0pt}S2c\_\hspace{0pt}lift/\hspace{0pt}stretch 输出} & 按干预分组;\allowbreak  计数+范围;\allowbreak  n=\allowbreak 4/\allowbreak 8 & 已核\\
C99 & 有理族四档（$\tau$=.06）：$\beta$=1 中位 1.66e-2 [2.87e-3;.327] 好 8；$\beta$=3 1.53e-2 [6.56e-3;.150] 好 9；$\beta$=8 5.56e-2 [1.80e-2;.375] 好 6；$\beta$=15 5.35e-2 [2.90e-2;.105] 好 8 \newline {\scriptsize §7.3 tab:d1beta} & {\scriptsize\texttt{aggregated/\hspace{0pt}v8\_\hspace{0pt}time\_\hspace{0pt}varying\_\hspace{0pt}d1r/\hspace{0pt}d1r\_\hspace{0pt}runs.csv}} & $\beta$ 档$\times$时空全域相对 L2;\allowbreak  中位+范围+计数;\allowbreak  n=\allowbreak 11/\allowbreak 档 & 已核\\
C100 & 配对 Wilcoxon $\beta$=1 对 3/8/15：p=.41/.32/.37 均不显著 \newline {\scriptsize §7.3 (ii)} & {\scriptsize\texttt{d1r\_\hspace{0pt}runs.csv}} & 同种子配对;\allowbreak  Wilcoxon 双侧;\allowbreak  n=\allowbreak 11 & 已核\\
C101 & 逐种子最优 $\beta$=3 拿 6/11 $\beta$=1 拿 4/11 $\beta$=15 拿 1/11；seed 30 由 .327 降至 .029；$\beta$=8/$\beta$=1 比值中位 3.95 最差 97 倍 \newline {\scriptsize §7.3 (ii)} & {\scriptsize\texttt{d1r\_\hspace{0pt}runs.csv}} & 逐种子最小值归属+配对比值;\allowbreak  计数+中位+最大;\allowbreak  n=\allowbreak 11 & 已核\\
C102 & 机制A：cos 中位 .9999 [.999;1.000]；总范数比 7.75 [7.15;7.84]；行范数比 2.90$\to$4.20 \newline {\scriptsize §7.3 (iii)} & {\scriptsize\texttt{aggregated/\hspace{0pt}v8\_\hspace{0pt}time\_\hspace{0pt}varying\_\hspace{0pt}d1r/\hspace{0pt}d1r\_\hspace{0pt}mechanism\_\hspace{0pt}a.csv}} & Gate 冻结 $\beta$=\allowbreak 1$\to$8 探针;\allowbreak  中位+范围;\allowbreak  n=\allowbreak 22 探针对 & 已核\\
C103 & 阈值敏感性：$\tau$=.05/.06/.10 下 $\beta$=8 为 4/6/8、$\beta$=15 为 5/8/10；$\tau$=.06 取在自然间隙（四档排序值 .056--.066 有间隙） \newline {\scriptsize §7.3 表注} & {\scriptsize\texttt{d1r\_\hspace{0pt}runs.csv}} & $\tau$ 三档计数+排序间隙;\allowbreak  计数;\allowbreak  n=\allowbreak 11/\allowbreak 档 & 已核\\
C104 & 空间线性族对照：$\beta$=8/15 零好盆地（增权有害）；早期另一 $\nu$/种子口径的扫描曾得增权有益 \newline {\scriptsize §7.3 (ii)} & {\scriptsize aggregated/\hspace{0pt}v8\_\hspace{0pt}time\_\hspace{0pt}varying/\hspace{0pt}summary.csv 与 summary\_\hspace{0pt}1/\hspace{0pt}\_\hspace{0pt}2/\hspace{0pt}\_\hspace{0pt}3.csv} & $\beta$ 档$\times$L2\_\allowbreak relative;\allowbreak  计数;\allowbreak  n=\allowbreak 11/\allowbreak 档 & 已核\\
C105 & 各组终态 L2 中位：对照 2.57、down\_3 2.75、down\_8 2.86、up\_3 2.44、up\_8 2.44；clip 与对照逐位相同 \newline {\scriptsize §7.3 实验 7.6} & {\scriptsize\texttt{aggregated/\hspace{0pt}v8\_\hspace{0pt}time\_\hspace{0pt}varying/\hspace{0pt}summary\_\hspace{0pt}11--\_\hspace{0pt}18.csv（seeds 200--210）}} & 逐组终态 L2;\allowbreak  中位;\allowbreak  n=\allowbreak 11/\allowbreak 组 & 已核\\
C106 & 三档中位 .0506/.0503/.0485 范围 [.0461;.0628]/[.0418;.0582]/[.0414;.0583] \newline {\scriptsize §7.4 tab:c45} & {\scriptsize\texttt{aggregated/\hspace{0pt}v8\_\hspace{0pt}burgers\_\hspace{0pt}2d/\hspace{0pt}summary\_\hspace{0pt}2.json}} & $\beta$ 档$\times$L2\_\allowbreak total;\allowbreak  中位+范围;\allowbreak  n=\allowbreak 11/\allowbreak 档 & 已核\\
C107 & 全部 33 run 进好盆地（表列 11/11/11）；$\beta$=1$\to$8 约 4\% 同量级波动 \newline {\scriptsize §7.4 结论} & {\scriptsize\texttt{summary\_\hspace{0pt}2.json}} & L2\_\allowbreak total 与 $\tau$=\allowbreak .06 比较;\allowbreak  计数;\allowbreak  n=\allowbreak 33 & 已核\\
C108 & 实验 4.2 线性 gap 口径 Jacobi 填补 >85\%；rational 族内最优恒 $\beta$=1 \newline {\scriptsize §8.3} & {\scriptsize\texttt{aggregated/\hspace{0pt}exp\_\hspace{0pt}V2\_\hspace{0pt}P1\_\hspace{0pt}constrained/\hspace{0pt}p1\_\hspace{0pt}constrained.csv}} & 同 C26--C29;\allowbreak  导出;\allowbreak  n=\allowbreak 4 & 已核\\
C109 & uniform .613 最高；down\_lin .388 Holm p=.006；down\_inv/rational 亦低于 uniform \newline {\scriptsize §8.4 (4)} & {\scriptsize\texttt{aggregated/\hspace{0pt}v7\_\hspace{0pt}v1\_\hspace{0pt}exp\_\hspace{0pt}G1b\_\hspace{0pt}pentry/\hspace{0pt}g1b\_\hspace{0pt}pentry\_\hspace{0pt}agg.csv}} & 同 C04--C08;\allowbreak  原值;\allowbreak  n=\allowbreak 80 & 已核\\
C110 & 实验 5.9 固定步长 L-BFGS 四族终态差异不显著 \newline {\scriptsize §8.4 (4)} & {\scriptsize\texttt{aggregated/\hspace{0pt}v7\_\hspace{0pt}v2\_\hspace{0pt}exp\_\hspace{0pt}V2\_\hspace{0pt}GstabB\_\hspace{0pt}stabilizer}} & 同 C53/\allowbreak C54;\allowbreak  配对 Wilcoxon;\allowbreak  n=\allowbreak 8 配对 & 已核\\
C111 & 无新定量声明（条件清单与结论汇总） \newline {\scriptsize §8.1/§8.2/§8.5/总结段} & — & — & 已核\\
C112 & 实验 7.5 机制A探针 N=2048 p=7851；配点为随机（非固定） \newline {\scriptsize 附录 B N/p 表} & {\scriptsize\texttt{code/\hspace{0pt}v8\_\hspace{0pt}2d/\hspace{0pt}experiments/\hspace{0pt}time\_\hspace{0pt}varying/\hspace{0pt}exp\_\hspace{0pt}D1\_\hspace{0pt}rational\_\hspace{0pt}beta\_\hspace{0pt}scan.py L161-182}} & 代码审查;\allowbreak  代码审查;\allowbreak  n=\allowbreak 22 探针对 & 已核\\
C113 & bc\_amp 审计：$\delta$($\nu$)=1-tanh(1/2$\nu$) 四档数值；实验 7.8 作废 264 runs \newline {\scriptsize 附录 B.5} & {\scriptsize\texttt{aggregated/\hspace{0pt}v7\_\hspace{0pt}v2\_\hspace{0pt}exp\_\hspace{0pt}V2\_\hspace{0pt}S1\_\hspace{0pt}sigma\_\hspace{0pt}nu\_\hspace{0pt}scan/\hspace{0pt}sigmanu\_\hspace{0pt}fixed\_\hspace{0pt}runs\_\hspace{0pt}invalid\_\hspace{0pt}bcamp.csv}} & 数值复算+行数;\allowbreak  复算+计数;\allowbreak  n=\allowbreak 4/\allowbreak 264 & 已核\\
C114 & 实验 4.3/5.1/5.2/5.7/7.8 样本量 \newline {\scriptsize 附录 实验图谱} & {\scriptsize\texttt{r1\_\hspace{0pt}runs.csv 330、g1c\_\hspace{0pt}clip\_\hspace{0pt}runs.csv 160、sigmanu\_\hspace{0pt}fixed\_\hspace{0pt}runs.csv 990、gstab\_\hspace{0pt}early\_\hspace{0pt}runs.csv 264（88/\hspace{0pt}$\nu$）}} & 各聚合 csv 行数与分组;\allowbreak  计数 & 已核\\
C115 & 实验 5.5 297、实验 5.6 1584、实验 5.3/6.1 各 80 的格数算术 \newline {\scriptsize 附录 app:tables} & {\scriptsize 设计网格} & 算术 & 已核\\
\end{longtable}
}
